\documentclass[../../main.tex]{subfiles}

\begin{document}

\section{拓扑的基}\label{section:ch02-拓扑的基}

上节的每一个例子都是用开集族 $\mathcal{T}$ 来刻画拓扑的，这样做一般显得不太方便。更多的时候是首先指定 $X$ 的子集的一个较小的族，然后用它来确定拓扑。

\begin{definition}[基与由基生成的拓扑]\label{definition:ch02-拓扑的基}
如果 $X$ 是一个集合，$X$ 的某拓扑的一个基（basis）是 $X$ 的子集的一个族 $\mathcal{B}$（其成员称为基元素（basis element）），满足条件：
\begin{enumerate}[(1)]
\item 对于每一个 $x \in X$，至少存在一个包含 $x$ 的基元素 $B$.
\item 若 $x$ 属于两个基元素 $B_{1}$ 和 $B_{2}$ 的交，则存在包含 $x$ 的一个基元素 $B_{3}$，使得 $B_{3} \subset B_{1} \cap B_{2}$.
\end{enumerate}

如果 $\mathcal{B}$ 满足以上两个条件，我们定义由 $\mathcal{B}$ 生成的拓扑 $\mathcal{T}$（topology $\mathcal{T}$ generated by $\mathcal{B}$）如下：如果对于每一个 $x \in U$，存在一个基元素 $B \in \mathcal{B}$，使得 $x \in B$ 并且 $B \subset U$，那么 $X$ 的子集 $U$ 称为 $X$ 的开集（即是 $\mathcal{T}$ 的一个元素）。
\end{definition}

我们将简短地验证 $\mathcal{T}$ 确实是 $X$ 的一个拓扑。在此之前，先看几个例子。

\par
\begin{figure}[htbp]
\centering
\includegraphics[width=0.4\textwidth]{ch02-fig13-1}
\caption{圆形域族满足基条件}\label{figure:ch02-圆域基}
\end{figure}

\begin{example}\label{example:ch02-圆域基}
设 $\mathcal{B}$ 为平面上所有圆形域（圆周的内部）所组成的族，则 $\mathcal{B}$ 满足基的定义中的两个条件。第二个条件如\cref{figure:ch02-圆域基}所示。在由 $\mathcal{B}$ 生成的拓扑中，平面的一个子集 $U$ 是开集，是指对于 $U$ 中任意 $x$，都有含于 $U$ 中的某一个圆域。

\end{example}

\par
\begin{figure}[htbp]
\centering
\includegraphics[width=0.4\textwidth]{ch02-fig13-2}
\caption{矩形域族满足基条件}\label{figure:ch02-矩形域基}
\end{figure}

\begin{example}\label{example:ch02-矩形域基}
设 $\mathcal{B}'$ 为平面上所有矩形域（矩形的内部）的族，其中矩形的边平行于两个坐标轴，则 $\mathcal{B}'$ 满足基的定义中的两个条件。第二个条件如\cref{figure:ch02-矩形域基}所示。这个条件是显然满足的，这是由于任何两个基元素的交，其本身就是一个基元素（或者是空集）。以后我们将会看到，在平面上由基 $\mathcal{B}'$ 生成的拓扑与\cref{example:ch02-圆域基}中基 $\mathcal{B}$ 生成的拓扑是一样的。

\end{example}

\begin{example}\label{example:ch02-单点基}
设 $X$ 是任意的一个集合。$X$ 的所有单点子集的族是 $X$ 上离散拓扑的一个基。
\end{example}

现在我们来验证：由基 $\mathcal{B}$ 生成的族 $\mathcal{T}$ 的确是 $X$ 的一个拓扑。若 $U$ 为空集，它显然满足开集的定义。同样，因为对于每一个 $x \in X$，存在包含 $x$ 的某一个基元素 $B$ 且 $B \subset X$，所以 $X$ 在 $\mathcal{T}$ 中。现在取 $\mathcal{T}$ 中元素的加标族 $\{U_{\alpha}\}_{\alpha \in J}$，并且证明
\begin{align*}
U = \bigcup_{\alpha \in J} U_{\alpha}
\end{align*}
属于 $\mathcal{T}$。对于 $x \in U$，存在一个指标 $\alpha$，使得 $x \in U_{\alpha}$。因为 $U_{\alpha}$ 是开集，所以存在一个基元素 $B$，使得 $x \in B \subset U_{\alpha}$。由于 $x \in B$ 并且 $B \subset U$，根据定义可见 $U$ 为开集。

任意选取 $\mathcal{T}$ 的两个元素 $U_{1}$ 和 $U_{2}$，以下证明 $U_{1} \cap U_{2}$ 属于 $\mathcal{T}$。给定 $x \in U_{1} \cap U_{2}$，选取一个包含 $x$ 的基元素 $B_{1}$，使得 $B_{1} \subset U_{1}$。再选取包含 $x$ 的基元素 $B_{2}$，使得 $B_{2} \subset U_{2}$。根据基所满足的第二个条件，存在一个包含 $x$ 的基元素 $B_{3}$，使得 $B_{3} \subset B_{1} \cap B_{2}$（参见\cref{figure:ch02-基交}）。于是 $x \in B_{3}$ 并且 $B_{3} \subset U_{1} \cap U_{2}$，从而根据定义可见 $U_{1} \cap U_{2}$ 属于 $\mathcal{T}$。

\par
\begin{figure}[htbp]
\centering
\includegraphics[width=0.55\textwidth]{ch02-fig13-3}
\caption{基元素的交仍含于开集}\label{figure:ch02-基交}
\end{figure}

最后，我们归纳地证明 $\mathcal{T}$ 中元素的任何有限交 $U_{1} \cap \cdots \cap U_{n}$ 在 $\mathcal{T}$ 中。当 $n = 1$ 时这是显然的。现在假定对于 $n-1$ 结论已经成立，证明对 $n$ 这结论还是正确的。因为
\begin{align*}
(U_{1} \cap \cdots \cap U_{n}) = (U_{1} \cap \cdots \cap U_{n-1}) \cap U_{n},
\end{align*}
由归纳假设，$U_{1} \cap \cdots \cap U_{n-1}$ 属于 $\mathcal{T}$，应用刚才证明的结果，$U_{1} \cap \cdots \cap U_{n-1}$ 与 $U_{n}$ 的交也属于 $\mathcal{T}$。

这样一来，我们就验证了由基 $\mathcal{B}$ 生成的开集族的确是一个拓扑。

由基生成拓扑的另一种描述方法是由以下引理给出的。

\begin{lemma}[基生成拓扑即基元素之并]\label{lemma:ch02-基生成并}
设 $X$ 是一个集合，$\mathcal{B}$ 是 $X$ 的拓扑 $\mathcal{T}$ 的一个基。则 $\mathcal{T}$ 等于 $\mathcal{B}$ 中元素所有并的族。
\end{lemma}

\begin{proof}
给定 $\mathcal{B}$ 中元素的一个族，这些 $\mathcal{B}$ 的元素也是 $\mathcal{T}$ 的元素。由于 $\mathcal{T}$ 是一个拓扑，于是它们的并也在 $\mathcal{T}$ 中。反之，给定 $U \in \mathcal{T}$，对于每一个 $x \in U$，选取 $\mathcal{B}$ 的一个元素 $B_{x}$，使得 $x \in B_{x} \subset U$。从而，$U$ 等于 $\mathcal{B}$ 中元素的一个并。
\end{proof}

上述引理表明 $X$ 中的每一开集 $U$ 都可以表示成某些基元素之并。然而，$U$ 这种表示不是唯一的。因而，拓扑学中“基”这个概念有别于它在代数学中的用法，在代数学中，将一个给定向量写成基向量的线性组合时其表达式是唯一的。

我们已经使用两种不同的方法描述了如何由基生成拓扑。有的时候我们需要反过来做，即由一个拓扑找出生成它的基。下面介绍如何由给定的拓扑得到基的一种方法，以后要多次用到。

\begin{lemma}[开集族为基]\label{lemma:ch02-开集族为基}
设 $X$ 是一个拓扑空间。$\mathcal{C}$ 是 $X$ 的开集的一个族，它满足对于 $X$ 的每一个开集 $U$ 及每一个 $x \in U$，存在 $\mathcal{C}$ 的一个元素 $C$，使得 $x \in C \subset U$。那么 $\mathcal{C}$ 就是 $X$ 上这个拓扑的一个基。
\end{lemma}

\begin{proof}
我们必须证明 $\mathcal{C}$ 是一个基。基的定义中的第一个条件是显然满足的：对于 $x \in X$，因为 $X$ 本身是开集，根据假设存在 $C \in \mathcal{C}$，使得 $x \in C \subset X$。为了验证第二个条件，设 $x \in C_{1} \cap C_{2}$，其中 $C_{1}$ 和 $C_{2}$ 是 $\mathcal{C}$ 的两个元素。因为 $C_{1}$ 和 $C_{2}$ 是开集，所以 $C_{1} \cap C_{2}$ 也是开集。于是根据假设，存在 $C_{3} \in \mathcal{C}$，使得 $x \in C_{3} \subset C_{1} \cap C_{2}$。设 $\mathcal{T}'$ 表示 $X$ 的由 $\mathcal{C}$ 所生成的拓扑，$\mathcal{T}$ 为 $X$ 的拓扑。上述引理表明 $\mathcal{T}'$ 细于 $\mathcal{T}$。反之，因为 $\mathcal{C}$ 的每一个元素都是 $\mathcal{T}$ 的一个元素，所以 $\mathcal{C}$ 的元素的任意并也在 $\mathcal{T}$ 中。于是根据\cref{lemma:ch02-基生成并} $\mathcal{T}' \subset \mathcal{T}$。这就证明了 $\mathcal{T}' = \mathcal{T}$。
\end{proof}

当拓扑是由基给出的时候，就可以用基作为判定拓扑粗细的一个标准。下面就是这样的一个准则。

\begin{lemma}[基的粗细比较]\label{lemma:ch02-基粗细}
设 $\mathcal{B}$ 和 $\mathcal{B}'$ 分别是 $X$ 的拓扑 $\mathcal{T}$ 和 $\mathcal{T}'$ 的基。则下列条件等价：
\begin{enumerate}[(1)]
\item $\mathcal{T}'$ 细于 $\mathcal{T}$.
\item 对于每一个 $x \in X$ 及包含 $x$ 的每一个基元素 $B \in \mathcal{B}$，存在一个基元素 $B' \in \mathcal{B}'$，使得 $x \in B' \subset B$.
\end{enumerate}
\end{lemma}

\begin{proof}
(2) $\Longrightarrow$ (1)。对于 $\mathcal{T}$ 的一个元素 $U$，我们证明 $U \in \mathcal{T}'$。取 $x \in U$。因为 $\mathcal{B}$ 生成 $\mathcal{T}$，则存在一个元素 $B \in \mathcal{B}$，使得 $x \in B \subset U$。条件 (2) 告诉我们，存在一个元素 $B' \in \mathcal{B}'$，使得 $x \in B' \subset B$。于是 $x \in B' \subset U$。根据定义 $U \in \mathcal{T}'$。

(1) $\Longrightarrow$ (2)。给定 $x \in X$ 和 $B \in \mathcal{B}$，其中 $x \in B$。根据定义，$B \in \mathcal{T}$，再根据条件 (1)，$\mathcal{T} \subset \mathcal{T}'$。因此 $B \in \mathcal{T}'$。因为 $\mathcal{B}'$ 生成 $\mathcal{T}'$，所以存在一个元素 $B' \in \mathcal{B}'$，使得 $x \in B' \subset B$。
\end{proof}

有的学生认为这个条件不容易记住。他们问道：“怎样才能记住这个包含关系呢？”回忆一下我们把拓扑空间看成为一辆装满石子的汽车就好记了。每一个石子看成拓扑的基元素，然后把每一个石子敲碎，将每颗碎粒看成新拓扑的一个基元素。新拓扑就比原拓扑细，而每一颗碎粒必能被原先的一个石子包含，这正是前面这个准则中所讲的事情。

\par
\begin{figure}[htbp]
\centering
\includegraphics[width=0.5\textwidth]{ch02-fig13-4}
\caption{圆域与矩形域生成同一拓扑}\label{figure:ch02-圆与矩形}
\end{figure}

\begin{example}\label{example:ch02-圆与矩形同拓扑}
现在可以看出，平面上所有圆形域的族 $\mathcal{B}$ 与所有矩形域的族 $\mathcal{B}'$，生成的是同一拓扑，\cref{figure:ch02-圆与矩形}给出了证明的图示。在学习了度量空间之后，我们还将对这个例子更正式地加以讨论。

\end{example}

下面我们来定义实直线 $\mathbb{R}$ 上三种有趣的拓扑。

\begin{definition}[标准拓扑、下限拓扑与 $K$-拓扑]\label{definition:ch02-实数三种拓扑}
设 $\mathcal{B}$ 为实直线上所有开区间
\begin{align*}
(a, b) = \{x \mid a < x < b\}
\end{align*}
的族。由 $\mathcal{B}$ 生成的拓扑称为实直线上的标准拓扑（standard topology）。当我们考虑 $\mathbb{R}$ 时，若无特别声明，我们将总假设考虑的是这个标准拓扑。设 $\mathcal{B}'$ 为所有半开区间
\begin{align*}
[a, b) = \{x \mid a \leqslant x < b\}
\end{align*}
的族，其中 $a < b$。由 $\mathcal{B}'$ 生成的拓扑称为 $\mathbb{R}$ 的下限拓扑（lower limit topology）。具有下限拓扑的 $\mathbb{R}$ 记作 $\mathbb{R}_{\ell}$。对于 $n \in \mathbb{Z}_{+}$，令 $K$ 表示所有形如 $\frac{1}{n}$ 的数所组成的集合，令 $\mathcal{B}''$ 是由所有开区间 $(a, b)$ 及形如 $(a, b) - K$ 形式的集合的族。由 $\mathcal{B}''$ 所生成的拓扑称为 $\mathbb{R}$ 上的 $K$-拓扑（$K$-topology）。具有 $K$-拓扑的 $\mathbb{R}$ 记作 $\mathbb{R}_{K}$。
\end{definition}

易见，这三个集族都是基，在每一集族中的两个基元素的交或者是另一个基元素，或者是空集。这些拓扑间的关系如下：

\begin{lemma}[下限拓扑与 $K$-拓扑严格细于标准拓扑]\label{lemma:ch02-下限K严格细}
$\mathbb{R}_{\ell}$ 和 $\mathbb{R}_{K}$ 的拓扑都严格细于标准拓扑，但它们之间不可比较。
\end{lemma}

\begin{proof}
以 $\mathcal{T}$，$\mathcal{T}'$，$\mathcal{T}''$ 分别表示 $\mathbb{R}$，$\mathbb{R}_{\ell}$，$\mathbb{R}_{K}$ 中的拓扑。任意给定 $\mathcal{T}$ 中的一个基元素 $(a, b)$ 以及一个点 $x \in (a, b)$，$\mathcal{T}'$ 中的元素 $[x, b)$ 包含着 $x$ 并且包含于 $(a, b)$。另一方面，任意给定 $\mathcal{T}'$ 中的一个基元素 $[x, d)$，却不存在 $\mathcal{T}$ 中的元素 $(a, b)$，使得它包含于 $[x, d)$ 同时又包含着 $x$。因此 $\mathcal{T}'$ 严格细于 $\mathcal{T}$。

对 $\mathbb{R}_{K}$ 进行类似的讨论。任意给定 $\mathcal{T}$ 中的一个基元素 $(a, b)$ 以及一个点 $x \in (a, b)$，则 $(a, b)$ 是 $\mathcal{T}''$ 中的一个基元素并且包含着 $x$。另一方面，任意给定 $\mathcal{T}''$ 中的一个基元素 $B = (-1, 1) - K$ 以及 $B$ 中一个点 $0$，却不存在包含于 $B$ 中并且包含着 $0$ 的开区间。

请读者自行验证 $\mathbb{R}_{\ell}$ 的拓扑与 $\mathbb{R}_{K}$ 的拓扑不可比较。
\end{proof}

这里有一个问题。对于由一个基 $\mathcal{B}$ 生成的拓扑，可以描述成 $\mathcal{B}$ 中元素任意并的族。那么，如果对于集一个给定的族，就像取任意并那样取它们的有限交将得到些什么呢？这个问题引出拓扑的子基这个概念。

\begin{definition}[子基与由子基生成的拓扑]\label{definition:ch02-子基}
$X$ 的某拓扑的一个子基（subbasis）$\mathcal{S}$ 是 $X$ 的子集的一个族，它的并等于 $X$。由子基 $\mathcal{S}$ 生成的拓扑（topology generated by the subbasis $\mathcal{S}$）$\mathcal{T}$ 定义为 $\mathcal{S}$ 中元素的有限交的所有并的族。
\end{definition}

当然必须验证 $\mathcal{T}$ 确实是一个拓扑。为此，只要证明 $\mathcal{S}$ 中元素的所有有限交的族 $\mathcal{B}$ 是一个基就行了，再根据\cref{lemma:ch02-基生成并}可见 $\mathcal{B}$ 中元素任意并的族 $\mathcal{T}$ 是一个拓扑。给定 $x \in X$，它一定属于 $\mathcal{S}$ 的一个元素，因而属于 $\mathcal{B}$ 的一个元素，这就是基的第一个条件。下面验证第二个条件。设
\begin{align*}
B_{1} = S_{1} \cap \cdots \cap S_{m} \quad \text{和} \quad B_{2} = S_{1}' \cap \cdots \cap S_{n}'
\end{align*}
为 $\mathcal{B}$ 的两个元素，它们的交
\begin{align*}
B_{1} \cap B_{2} = (S_{1} \cap \cdots \cap S_{m}) \cap (S_{1}' \cap \cdots \cap S_{n}')
\end{align*}
仍然是 $\mathcal{S}$ 中元素的一个有限交，因而属于 $\mathcal{B}$。

\subsection*{习题}

\begin{exercise}\label{exercise:ch02-拓扑的基-1}
设 $X$ 是一个拓扑空间，$A$ 是 $X$ 的一个子集。假定对于每一个 $x \in A$，存在包含 $x$ 的一个开集 $U$，使得 $U \subset A$。证明 $A$ 是 $X$ 中的一个开集。
\end{exercise}

\begin{exercise}\label{exercise:ch02-拓扑的基-2}
将第\cref{section:ch02-拓扑空间}节例 1 中集合 $X = \{a, b, c\}$ 的 9 个拓扑加以比较。即对于每两个拓扑，判定它们是否可以比较。如果可以比较，请指出哪一个较细。
\end{exercise}

\begin{exercise}\label{exercise:ch02-拓扑的基-3}
证明第\cref{section:ch02-拓扑空间}节例 4 中给出的族 $\mathcal{T}_{c}$ 是集合 $X$ 的一个拓扑。族
\begin{align*}
\mathcal{T}_{\infty} = \{U \mid X - U \text{ 或为无限集，或为空集，或为 } X\}
\end{align*}
是 $X$ 上的一个拓扑吗？
\end{exercise}

\begin{exercise}\label{exercise:ch02-拓扑的基-4}
(a) 设 $\{\mathcal{T}_{\alpha}\}$ 是 $X$ 的拓扑的一个族。证明 $\bigcap \mathcal{T}_{\alpha}$ 是 $X$ 上的一个拓扑。$\bigcup \mathcal{T}_{\alpha}$ 是 $X$ 上的一个拓扑吗？

(b) 设 $\{\mathcal{T}_{\alpha}\}$ 是 $X$ 的拓扑的一个族。证明 $X$ 上存在包含所有 $\mathcal{T}_{\alpha}$ 的唯一一个最小的拓扑，也存在包含于所有 $\mathcal{T}_{\alpha}$ 的唯一一个最大的拓扑。

(c) 设 $X = \{a, b, c\}$，令
\begin{align*}
\mathcal{T}_{1} = \{\varnothing, X, \{a\}, \{a, b\}\} \quad \text{和} \quad \mathcal{T}_{2} = \{\varnothing, X, \{a\}, \{b, c\}\},
\end{align*}
求出包含着 $\mathcal{T}_{1}$ 和 $\mathcal{T}_{2}$ 的最小拓扑，以及包含于 $\mathcal{T}_{1}$ 和 $\mathcal{T}_{2}$ 的最大拓扑。
\end{exercise}

\begin{exercise}\label{exercise:ch02-拓扑的基-5}
证明：若 $\mathcal{A}$ 是 $X$ 的拓扑的一个基，则由 $\mathcal{A}$ 生成的拓扑等于包含着 $\mathcal{A}$ 的 $X$ 的所有拓扑的交。

若 $\mathcal{A}$ 是一个子基，请证明同样的结论。
\end{exercise}

\begin{exercise}\label{exercise:ch02-拓扑的基-6}
证明：$\mathbb{R}_{\ell}$ 的拓扑与 $\mathbb{R}_{K}$ 的拓扑不可比较。
\end{exercise}

\begin{exercise}\label{exercise:ch02-拓扑的基-7}
考虑 $\mathbb{R}$ 的下列拓扑：
\begin{align*}
\mathcal{T}_{1} = \text{标准拓扑},
\end{align*}
\begin{align*}
\mathcal{T}_{2} = \mathbb{R}_{K} \text{ 的拓扑},
\end{align*}
\begin{align*}
\mathcal{T}_{3} = \text{有限补拓扑},
\end{align*}
\begin{align*}
\mathcal{T}_{4} = \text{上限拓扑，即以全体形如 } (a, b] \text{ 的集合为基的拓扑},
\end{align*}
\begin{align*}
\mathcal{T}_{5} = \text{以全体形如 } (-\infty, a) = \{x \mid x < a\} \text{ 的集合为基的拓扑}.
\end{align*}
对于上述每一个拓扑，确定它包含着上述拓扑中哪些与它不同的拓扑。
\end{exercise}

\begin{exercise}\label{exercise:ch02-拓扑的基-8}
(a) 应用\cref{lemma:ch02-开集族为基}证明可数族
\begin{align*}
\mathcal{B} = \{(a, b) \mid a < b, \ a \text{ 和 } b \text{ 都是有理数}\}
\end{align*}
为生成 $\mathbb{R}$ 的标准拓扑的一个基。

(b) 证明集族
\begin{align*}
\mathcal{C} = \{[a, b) \mid a < b, \ a \text{ 和 } b \text{ 都是有理数}\}
\end{align*}
为 $\mathbb{R}$ 的某一个拓扑的基，并且所生成的拓扑与下限拓扑不同。
\end{exercise}

\end{document}